\documentclass[10pt,a4paper]{amsart}
\usepackage[margin=19mm]{geometry}
\usepackage{amsmath,amssymb,mathtools}
\usepackage[scr=rsfs,cal=euler]{mathalfa}
\usepackage{xfrac}
\usepackage[unicode,hidelinks,bookmarks=false]{hyperref}
\newcommand{\ie}{i.e.\ }
\newcommand{\cf}{cf.\ }
\newcommand{\resp}{respectively\ }
\newcommand{\Subsection}{\subsection}
\providecommand{\nobreakdash}{\nobreak\hbox{-}}
\setlength{\emergencystretch}{2em}
\multlinegap0pt
\usepackage[shortlabels]{enumitem}
\usepackage{amssymb}
\usepackage{mathtools}
\newtheorem{assumption}{Assumption}
\newtheorem{lemma}{Lemma}
\newcommand{\E}{\mathbb{E}}
\newcommand{\Hb}{\mathbb{H}}
\newcommand{\Pbb}{\mathbb{P}}
\newcommand{\R}{\mathbb{R}}
\newcommand{\dd}{\mathrm{d}}
\newcommand{\dt}{\,\dd t}
\title{Existence of densities and atoms for the running maximum of time-inhomogeneous jump diffusions}
\author{Takuya Nakagawa, Ryoichi Suzuki}
\date{Source: arXiv:2608.24294v1 (CC BY 4.0)}
\begin{document}
\maketitle

\noindent\emph{Source and licence.} Existence of densities and atoms for the running maximum of time-inhomogeneous jump diffusions. Version arXiv:2608.24294v1 (CC BY 4.0), distributed by arXiv under the Creative Commons Attribution 4.0 International licence, http://creativecommons.org/licenses/by/4.0/, as recorded in the arXiv licence metadata for this exact version and shown on the abstract page https://arxiv.org/abs/2608.24294v1. Full licence notice: \texttt{licenses/arxiv-2608.24294-CC-BY-4.0.txt}.\par\smallskip

\noindent\textit{Editorial scope.} Counted unit: the article's lemma lem:DXrep (source0.tex lines 1448--1467). The article's complete original proof of it is delivered with it (source0.tex lines 1469--1611, one \texttt{proof} environment of 143 lines). No mathematics is written, repaired or re-derived: the statement and the proof are the article's own lines, in the article's own notation, and no retained line is substituted or re-flowed.  Retained uncounted as the source-local context the proof needs: lines 97--99; lines 388--397; lines 618--620; lines 711--718.  Nothing was omitted from any retained span, so the package carries no omission record.  No retained line contains a citation, so the wrapper carries no bibliography and no bibliography entry.  No retained line contains a figure environment, an image-inclusion command or drawing code, so no image is delivered and the package declares no asset dependency.  Editorial formatting only, in addition: an AMS \texttt{amsart} preamble that restates the article's own declarations and macros used by the retained lines, namely the environments \texttt{assumption}, \texttt{lemma} on separate counters, together with the standard packages those lines need.  The retained environments print in this document as Assumption 1, Lemma 1 in order of appearance, whereas the article prints the same items with the numbering of its own full document.  The complete native source of the article as distributed in the arXiv e-print for this exact version is bundled unchanged as \texttt{sources/208462/source0.tex}.
\begin{equation}\label{eq:SDE-general}
\dd X_t=b(t,X_t)\dt+\sigma(t,X_t)\,\dd W_t+\kappa(t)\,\dd L_t,\qquad X_0=x\in\R.
\end{equation}

\begin{assumption}[Coefficients: elliptic Brownian regime]\label{ass:elliptic}
Assume the following. Here $\|\cdot\|_\infty$ denotes the supremum norm over the relevant domain.
\begin{enumerate}[label=(E\arabic*),leftmargin=2.7em]
\item $b,\sigma:[0,T]\times\R\to\R$ are measurable in $t$ and $C^1$ in $x$ with
$\|\partial_x b\|_\infty+\|\partial_x\sigma\|_\infty<\infty$ and $\|b(\cdot,0)\|_\infty+\|\sigma(\cdot,0)\|_\infty<\infty$;
\item (\emph{Uniform ellipticity}) there exists $\sigma_0>0$ such that $\inf_{(t,x)\in[0,T]\times\R}|\sigma(t,x)|\ge\sigma_0$;
\item (\emph{Boundedness}) $\|\sigma\|_\infty<\infty$;
\item $\kappa:[0,T]\to\R$ is a bounded \emph{deterministic} measurable function, i.e., $\|\kappa\|_\infty<\infty$.
\end{enumerate}
\end{assumption}

\begin{equation}\label{eq:J}
\dd J_t=\partial_x b(t,X_t)J_t\dt+\partial_x\sigma(t,X_t)J_t\,\dd W_t,\qquad J_0=1.
\end{equation}

\begin{equation}\label{eq:linearized}
\begin{aligned}
D_\Theta X_t
&=\int_0^t \partial_x b(s,X_s)\,D_\Theta X_s\,\dd s
+\int_0^t \partial_x\sigma(s,X_s)\,D_\Theta X_s\,\dd W_s\\
&\quad+\int_0^t \sigma(s,X_s)h(s)\,\dd s,\qquad t\in[0,T].
\end{aligned}
\end{equation}

\begin{lemma}[Directional derivative via an adapted shift]\label{lem:DXrep}
Assume that Assumption~\ref{ass:elliptic} holds and let $h\in \Hb_p$ for some $p>1$.
For $\varepsilon\in\R$ define $W_t^\varepsilon:=W_t+\varepsilon\int_0^t h(s)\,\dd s$ and let $X^\varepsilon$ be the unique strong
solution to
\begin{equation}\label{eq:SDE-eps}
X_t^\varepsilon=x+\int_0^t b(s,X_s^\varepsilon)\,\dd s+\int_0^t \sigma(s,X_s^\varepsilon)\,\dd W_s
+\varepsilon\int_0^t \sigma(s,X_s^\varepsilon)h(s)\,\dd s+\int_0^t \kappa(s)\,\dd L_s,\qquad t\in[0,T].
\end{equation}
Then, as $\varepsilon\to0$,
\[
\frac{X^\varepsilon-X}{\varepsilon}\longrightarrow Y
\quad\text{in }L^p\big(\Omega;C([0,T])\big),
\]
where $Y$ is the unique adapted solution to the linear SDE \eqref{eq:linearized}.
Moreover,
\[
Y_t=J_t\int_0^t J_s^{-1}\sigma(s,X_s)h(s)\,\dd s,\qquad t\in[0,T],
\]
where $J$ is the Jacobian defined by \eqref{eq:J}.
\end{lemma}

\begin{proof}
Let $\Lambda:=\|\partial_x b\|_\infty\vee \|\partial_x\sigma\|_\infty$.
Since $\partial_x b$ and $\partial_x\sigma$ are bounded, $b(t,\cdot)$ and $\sigma(t,\cdot)$ are globally Lipschitz with constant $\Lambda$.

\medskip
\noindent\emph{Step 1: first-order estimate for $X^\varepsilon-X$.}
Set $\Delta_t^\varepsilon:=X_t^\varepsilon-X_t$. Subtracting \eqref{eq:SDE-general} from \eqref{eq:SDE-eps} gives
\[
\Delta_t^\varepsilon=\int_0^t\!\big(b(s,X_s^\varepsilon)-b(s,X_s)\big)\,\dd s
+\int_0^t\!\big(\sigma(s,X_s^\varepsilon)-\sigma(s,X_s)\big)\,\dd W_s
+\varepsilon\int_0^t \sigma(s,X_s^\varepsilon)h(s)\,\dd s.
\]
Note that the jump terms cancel in the difference (the same L\'evy process $L$ drives both equations), so $\Delta^\varepsilon$ has \emph{continuous}
paths.
Using BDG, Lipschitz continuity, and the bound $|\sigma|\le\|\sigma\|_\infty$, we obtain for $p>1$ a constant $C=C(p,T,\Lambda,\|\sigma\|_\infty)$ such that
\begin{equation}\label{eq:Delta-eps-bound}
\E\Big[\sup_{0\le t\le T}|\Delta_t^\varepsilon|^p\Big]
\le C\int_0^T \E\Big[\sup_{0\le u\le s}|\Delta_u^\varepsilon|^p\Big]\,\dd s
+C|\varepsilon|^p\,\E\Big[\Big(\int_0^T |h(s)|^2\,\dd s\Big)^{p/2}\Big].
\end{equation}
By Gronwall's lemma,
\begin{equation}
\E\Big[\sup_{0\le t\le T}|\Delta_t^\varepsilon|^p\Big]\le C|\varepsilon|^p\|h\|_{\Hb_p}^p,
\end{equation}
so in particular $\Delta^\varepsilon\to0$ in $L^p(\Omega;C([0,T]))$ as $\varepsilon\to0$.

\medskip
\noindent\emph{Step 2: equation for the difference quotient.}
For $\varepsilon\neq 0$ define $Y_t^\varepsilon:=\Delta_t^\varepsilon/\varepsilon$. Define the predictable coefficients
\[
a_s^\varepsilon:=\int_0^1 \partial_x b\big(s, X_s+\theta\Delta_s^\varepsilon\big)\,\dd\theta,
\qquad
c_s^\varepsilon:=\int_0^1 \partial_x \sigma\big(s, X_s+\theta\Delta_s^\varepsilon\big)\,\dd\theta.
\]
Then
\[
\frac{b(s,X_s^\varepsilon)-b(s,X_s)}{\varepsilon}=a_s^\varepsilon\,Y_s^\varepsilon,
\qquad
\frac{\sigma(s,X_s^\varepsilon)-\sigma(s,X_s)}{\varepsilon}=c_s^\varepsilon\,Y_s^\varepsilon,
\]
and hence $Y^\varepsilon$ satisfies the linear SDE
\begin{equation}\label{eq:Yeps}
Y_t^\varepsilon=\int_0^t a_s^\varepsilon Y_s^\varepsilon\,\dd s+\int_0^t c_s^\varepsilon Y_s^\varepsilon\,\dd W_s+\int_0^t \sigma(s,X_s^\varepsilon)h(s)\,\dd s.
\end{equation}
In particular, $|a_s^\varepsilon|\vee |c_s^\varepsilon|\le \Lambda$.

\medskip
\noindent\emph{Step 3: convergence $Y^\varepsilon\to Y$.}
Let $Y$ denote the unique solution to \eqref{eq:linearized} and set $Z^\varepsilon:=Y^\varepsilon-Y$.
Subtracting \eqref{eq:linearized} from \eqref{eq:Yeps} yields
\[
Z_t^\varepsilon=\int_0^t a_s^\varepsilon Z_s^\varepsilon\,\dd s+\int_0^t c_s^\varepsilon Z_s^\varepsilon\,\dd W_s+R_t^\varepsilon,
\]
where
\[
R_t^\varepsilon:=\int_0^t (a_s^\varepsilon-a_s)Y_s\,\dd s+\int_0^t (c_s^\varepsilon-c_s)Y_s\,\dd W_s
+\int_0^t \big(\sigma(s,X_s^\varepsilon)-\sigma(s,X_s)\big)h(s)\,\dd s,
\]
and we write $a_s:=\partial_x b(s,X_s)$, $c_s:=\partial_x\sigma(s,X_s)$.
Using BDG and Gronwall as in \eqref{eq:Delta-eps-bound}, we obtain
\begin{equation}\label{eq:Zeps-gron}
\E\Big[\sup_{0\le t\le T}|Z_t^\varepsilon|^p\Big]\le C\,\E\Big[\sup_{0\le t\le T}|R_t^\varepsilon|^p\Big],
\end{equation}
with $C=C(p,T,\Lambda)<\infty$.

We now show $\E[\sup_{t\le T}|R_t^\varepsilon|^p]\to0$ as $\varepsilon\to0$ by treating the three terms in $R^\varepsilon$ in order.
For the first term, by H\"older's inequality,
\[
\sup_{0\le t\le T}\Big|\int_0^t (a_s^\varepsilon-a_s)Y_s\,\dd s\Big|
\le \int_0^T |a_s^\varepsilon-a_s|\,|Y_s|\,\dd s,
\]
so that
\begin{equation}\label{eq:R1}
\E\Big[\sup_{0\le t\le T}\Big|\int_0^t (a_s^\varepsilon-a_s)Y_s\,\dd s\Big|^p\Big]
\le T^{p-1}\int_0^T \E\big[|a_s^\varepsilon-a_s|^p\,|Y_s|^p\big]\,\dd s.
\end{equation}
Since $\Delta^\varepsilon\to0$ in probability uniformly on $[0,T]$, continuity and boundedness of $\partial_x b$ imply $a_s^\varepsilon\to a_s$ in probability for every $s$.
Moreover $|a_s^\varepsilon-a_s|\le 2\Lambda$.
We also note that $Y$ has finite $p$th moments in the supremum norm. Indeed, by BDG and $|a_s|\vee|c_s|\le \Lambda$,
\[
\E\Big[\sup_{0\le t\le T}|Y_t|^p\Big]
\le C\int_0^T \E\Big[\sup_{0\le u\le s}|Y_u|^p\Big]\,\dd s
+ C\,\E\Big[\Big(\int_0^T |\sigma(s,X_s)|^2|h(s)|^2\,\dd s\Big)^{p/2}\Big],
\]
for a constant $C=C(p,T,\Lambda)$.
Since $|\sigma|\le \|\sigma\|_\infty$ and $h\in \Hb_p$, the second term is bounded by $C\,\|\sigma\|_\infty^p\|h\|_{\Hb_p}^p<\infty$, and Gronwall's lemma yields $\E[\sup_{t\le T}|Y_t|^p]<\infty$.
Using $|Y_s|^p\le \sup_{0\le u\le T}|Y_u|^p$ and $\E[\sup_{u\le T}|Y_u|^p]<\infty$,
Vitali's theorem (uniform integrability) implies that the integrand in \eqref{eq:R1} converges to $0$ in $L^1$ for each $s$,
and is dominated by $(2\Lambda)^p\E[\sup_{u\le T}|Y_u|^p]$.
Hence the right-hand side of \eqref{eq:R1} tends to $0$ as $\varepsilon\to0$.

For the second term, the BDG inequality yields
\[
\E\Big[\sup_{0\le t\le T}\Big|\int_0^t (c_s^\varepsilon-c_s)Y_s\,\dd W_s\Big|^p\Big]
\le C_p\,\E\Big[\Big(\int_0^T |(c_s^\varepsilon-c_s)Y_s|^2\,\dd s\Big)^{p/2}\Big].
\]
Similarly, continuity and boundedness of $\partial_x\sigma$ yield $c_s^\varepsilon\to c_s$ in probability for every $s$, and $|c_s^\varepsilon-c_s|\le 2\Lambda$.
In particular, $c^\varepsilon-c\to0$ in measure on $[0,T]\times\Omega$ (with respect to $\dd s\otimes\Pbb$), hence along any sequence $\varepsilon_n\downarrow0$ one can extract a subsequence (not relabeled) such that $c_s^{\varepsilon_n}\to c_s$ for a.e.\ $(s,\omega)$.
Since for each fixed $\omega$ the path $s\mapsto Y_s(\omega)$ is continuous (hence bounded) on $[0,T]$, dominated convergence yields $\int_0^T |(c_s^{\varepsilon_n}-c_s)Y_s|^2\,\dd s\to0$ a.s.\ along the subsequence, and therefore the same integral converges to $0$ in probability as $\varepsilon\to0$.
Consequently $\big(\int_0^T |(c_s^\varepsilon-c_s)Y_s|^2\,\dd s\big)^{p/2}\to0$ in probability.
Moreover,
\[
\Big(\int_0^T |(c_s^\varepsilon-c_s)Y_s|^2\,\dd s\Big)^{p/2}
\le (4\Lambda^2)^{p/2}\Big(\int_0^T |Y_s|^2\,\dd s\Big)^{p/2}
\le (4\Lambda^2T)^{p/2}\sup_{0\le s\le T}|Y_s|^p,
\]
which is integrable. By Vitali's theorem (uniform integrability), we obtain
\[
\E\Big[\sup_{0\le t\le T}\Big|\int_0^t (c_s^\varepsilon-c_s)Y_s\,\dd W_s\Big|^p\Big]\longrightarrow 0.
\]

For the third term, the Cauchy--Schwarz inequality in $s$ gives
\[
\sup_{0\le t\le T}\Big|\int_0^t \big(\sigma(s,X_s^\varepsilon)-\sigma(s,X_s)\big)h(s)\,\dd s\Big|^p
\le \Big(\int_0^T |\sigma(s,X_s^\varepsilon)-\sigma(s,X_s)|^2\,\dd s\Big)^{p/2}\Big(\int_0^T |h(s)|^2\,\dd s\Big)^{p/2}.
\]
Under (E3), $|\sigma(s,X_s^\varepsilon)-\sigma(s,X_s)|\le 2\|\sigma\|_\infty$, hence
\[
\Big(\int_0^T |\sigma(s,X_s^\varepsilon)-\sigma(s,X_s)|^2\,\dd s\Big)^{p/2}\Big(\int_0^T |h(s)|^2\,\dd s\Big)^{p/2}
\le (4\|\sigma\|_\infty^2 T)^{p/2}\Big(\int_0^T |h(s)|^2\,\dd s\Big)^{p/2},
\]
where the right-hand side belongs to $L^1(\Omega)$ since $h\in \Hb_p$.
Moreover, since $\Delta^\varepsilon\to0$ in probability uniformly on $[0,T]$, Lipschitz continuity of $\sigma$
implies $\int_0^T |\sigma(s,X_s^\varepsilon)-\sigma(s,X_s)|^2\,\dd s\le \Lambda^2T\sup_{s\le T}|\Delta_s^\varepsilon|^2\to0$ in probability.
Therefore the product converges to $0$ in probability while being dominated by the $L^1(\Omega)$ random variable $(4\|\sigma\|_\infty^2 T)^{p/2}(\int_0^T |h(s)|^2\,\dd s)^{p/2}$.
By Vitali's theorem (uniform integrability implied by the $L^1$ domination), we obtain
\[
\E\Big[\sup_{0\le t\le T}\Big|\int_0^t \big(\sigma(s,X_s^\varepsilon)-\sigma(s,X_s)\big)h(s)\,\dd s\Big|^p\Big]\longrightarrow 0.
\]

Combining the three terms, we obtain $\E[\sup_{t\le T}|R_t^\varepsilon|^p]\to0$.
Hence, by \eqref{eq:Zeps-gron}, $\E[\sup_{t\le T}|Z_t^\varepsilon|^p]\to0$.
Thus $Y^\varepsilon\to Y$ in $L^p(\Omega;C([0,T]))$.

\medskip
\noindent\emph{Step 4: variation of constants.}
Let $J$ solve \eqref{eq:J}. Applying It\^o's formula to $J_t^{-1}Y_t$ and using \eqref{eq:J} and \eqref{eq:linearized} gives
$\dd(J_t^{-1}Y_t)=J_t^{-1}\sigma(t,X_t)h(t)\,\dd t$ and hence
\[
Y_t=J_t\int_0^t J_s^{-1}\sigma(s,X_s)h(s)\,\dd s.
\]
The proof is complete.
\end{proof}

\end{document}
